%------------------------
% Resume Template
% Author : Anubhav Singh
% Github : https://github.com/xprilion
% License : MIT
%------------------------

\documentclass[a4paper,20pt]{article}

\usepackage{latexsym}
\usepackage[empty]{fullpage}
\usepackage{titlesec}
\usepackage{marvosym}
\usepackage[usenames,dvipsnames]{color}
\usepackage{verbatim}
\usepackage{enumitem}
\usepackage{tabularx}
\usepackage[pdftex]{hyperref}
\usepackage{fancyhdr}
\usepackage[utf8]{inputenc}
\usepackage[T1]{fontenc}

\pagestyle{fancy}
\fancyhf{} % clear all header and footer fields
\fancyfoot{}
\renewcommand{\headrulewidth}{0pt}
\renewcommand{\footrulewidth}{0pt}

% Adjust margins
\addtolength{\oddsidemargin}{-0.530in}
\addtolength{\evensidemargin}{-0.375in}
\addtolength{\textwidth}{1in}
\addtolength{\topmargin}{-.45in}
\addtolength{\textheight}{1in}

\urlstyle{rm}

\raggedbottom
\raggedright
\setlength{\tabcolsep}{0in}

% Sections formatting
\titleformat{\section}{
  \vspace{-10pt}\bfseries\raggedright\large
}{}{0em}{}[\color{black}\titlerule \vspace{-6pt}]

%-------------------------
% Custom commands
\newcommand{\resumeItem}[2]{
  \item\small{
    \textbf{#1}{#2 \vspace{-2pt}}
  }
}

\newcommand{\resumeSubheading}[4]{
  \vspace{-1pt}\item
    \begin{tabular*}{0.97\textwidth}{l@{\extracolsep{\fill}}r}
      \textbf{#1} & #2 \\
      \textsc{#3} & \textit{#4} \\
    \end{tabular*}\vspace{-5pt}
}


\newcommand{\resumeSubItem}[2]{\resumeItem{#1}{#2}\vspace{-3pt}}

% Per-entry skills line (no bullet, small, bold label)
\newcommand{\resumeSkills}[1]{\item[]\small{\textbf{Skills:}~#1}\vspace{-4pt}}

% Publication entry: year column + bold title + venue (reference style)
\newcommand{\resumePub}[3]{%
  \noindent\begin{tabularx}{\textwidth}{@{}p{2.4em}@{\hspace{0.4em}}X@{}}
    {\small\textcolor[gray]{0.4}{#1}} & {\small\textbf{#2}, #3}
  \end{tabularx}\vspace{1pt}
}

% Conference entry: year column + conference, duration + location (right)
\newcommand{\resumeConf}[4]{%
  \noindent\begin{tabularx}{\textwidth}{@{}p{2.4em}@{\hspace{0.4em}}X@{}r@{}}
    {\small\textcolor[gray]{0.4}{#1}} & {\small #2, #3} & {\small\textit{#4}}
  \end{tabularx}\vspace{1pt}
}

\renewcommand{\labelitemii}{$\circ$}

\newcommand{\resumeSubHeadingListStart}{\begin{itemize}[leftmargin=*]}
\newcommand{\resumeSubHeadingListEnd}{\end{itemize}}
\newcommand{\resumeItemListStart}{\begin{itemize}}
\newcommand{\resumeItemListEnd}{\end{itemize}\vspace{-7pt}}

%-----------------------------
%%%%%%  CV STARTS HERE  %%%%%%


\begin{document}

%----------HEADING-----------------
\begin{tabular*}{\textwidth}{l@{\extracolsep{\fill}}r}
  \textbf{{\LARGE Arthur Cancellieri Pires}} & E-mail: \href{mailto:arthur-pires@live.com}{arthur-pires@live.com}\\
  \textit{Data Scientist \textbar{} PhD Candidate} & \href{https://www.linkedin.com/in/arthur-cancellieri-pires-730963160}{LinkedIn: Arthur Cancellieri Pires}\\
  \href{https://github.com/pirao}{GitHub: @Pirao} & \href{https://github.com/pirao/AnomalyDetection2026}{Public project: Industrial Sensor Anomaly Detection API}\\
\end{tabular*}


%----------Work Experience-----------------

\section{Professional Experience}
\resumeSubHeadingListStart
  \resumeSubheading{Railway Laboratory (Lafer) -- partnership with VALE}{On-site}
  {Lead Data Scientist}{January 2022 to Present}
  \resumeItemListStart

    \resumeItem{}{Led a team of 12 delivering a predictive maintenance program for VALE, one of the world's largest mining and logistics companies, converting stakeholder needs into shipped ML systems.}

    \resumeItem{}{Cut equipment failure detection latency from 3 months to real time at 91\% precision using deep learning models and automated data quality checks on multivariate sensor streams.}

    \resumeItem{}{Flagged critical defects weeks before scheduled maintenance, including cases absent from the maintenance database, with unsupervised Transformer autoencoders on time series data.}

  \resumeItemListEnd
  \resumeSkills{Python | PyTorch | PyTorch Lightning | scikit-learn | Deep Learning | Transformers | Autoencoders | Anomaly Detection | Time Series | Deployment | FastAPI | Docker}

\vspace{-2pt}

\resumeSubheading{Railway Laboratory (Lafer) -- partnership with VALE}{On-site}
{Research Data Scientist}{June 2020 to January 2022}
\resumeItemListStart

  \resumeItem{}{Built a daily ETL pipeline processing 20,000-40,000 sensor records per asset per day, integrating multiple external sources into one analytical dataset that powered all downstream models.}
  
  \resumeItem{}{Reduced reliance on costly specialized inspections by predicting asset condition at an R2 of about 0.98 with deep learning models on sensor data.}

\resumeItemListEnd
\resumeSkills{Python | scikit-learn | PyTorch | SQL | PostgreSQL | ETL | Deep Learning | Signal Processing}

\vspace{-2pt}

\resumeSubheading{Tribology and Railway Dynamics Laboratory (LabTDF) -- partnership with VALE}{On-site}
{Machine Learning Intern}{Jan 2019 to Jun 2020}
\resumeItemListStart

  \resumeItem{}{Estimated hard-to-measure contact forces at an R2 of about 0.97 from onboard sensors using a simulation-validated ML pipeline, enabling VALE's predictive maintenance program on instrumented wagons.}

  \resumeItem{}{Optimized wheel/rail wear profiles with genetic algorithms, cutting wear 20\% and improving a key safety ratio ~35\% without sacrificing fatigue performance.}

  \resumeItem{}{Proved the case for a maintenance threshold change through wear/fatigue analysis, driving adoption into operations and extending service life 29\%-50\% (29\% from the threshold change alone, up to 50\% when combined with the optimized wheel profile)}

\resumeItemListEnd
\resumeSkills{Python | MATLAB | Machine Learning | scikit-learn | Optuna | Genetic Algorithms | Optimization | Multibody Simulation}
\resumeSubHeadingListEnd

\vspace{-3pt}

\vspace{-5pt}

%-----------EDUCATION-----------------
\section{Education}
  \resumeSubHeadingListStart
    \resumeSubheading
      {University of Campinas (UNICAMP)}{Campinas, Brazil}
      {Ph.D. in Mechanical Engineering}{Expected: August 2026}
      \resumeItemListStart
        \resumeItem{Thesis: }{A methodology for detecting track defects from instrumented railway vehicle data.}
      \resumeItemListEnd
    \resumeSubheading
      {University of Campinas (UNICAMP)}{Campinas, Brazil}
      {M.Sc. in Mechanical Engineering}{Jan 2022}
      \resumeItemListStart
        \resumeItem{Thesis: }{Measuring railway track irregularities from instrumented railway vehicle data using machine learning techniques.}
      \resumeItemListEnd
    \vspace{-4pt}
    \resumeSubheading
      {Federal University of Espírito Santo (UFES)}{Vitória, Brazil}
      {B.Sc. in Mechanical Engineering}{May 2020}
      \resumeItemListStart
        \resumeItem{Thesis: }{Indirect identification of wheel--rail contact forces of an instrumented heavy haul railway vehicle using machine learning.}
      \resumeItemListEnd
    \resumeSubHeadingListEnd

\vspace{-5pt}

%-----------LANGUAGES-----------------
\section{Languages}
\vspace{2pt}
{\small Portuguese (Native), English (Fluent), Italian (Basic)}

\vspace{1pt}

%-----------PUBLICATIONS-----------------
\section{Publications \& Conferences}
\vspace{2pt}

{\scshape Selected Publications}
\vspace{3pt}

\resumePub{2024}{Measuring vertical track irregularities from instrumented heavy haul railway vehicle data using machine learning}{Engineering Applications of Artificial Intelligence}
\resumePub{2023}{Influence of wheel tread wear on Rolling Contact Fatigue and on the dynamics of railway vehicles}{Wear}
\resumePub{2022}{Nonintrusive Operator Inference for Chaotic Systems}{IEEE Transactions on Artificial Intelligence}
\resumePub{2022}{Generation of realistic artificial track irregularities for multibody simulation using measured geometric data: from mid-chord to space-curve}{In Proceedings of the Fifth International Conference on Railway Technology}
\resumePub{2021}{Indirect identification of wheel--rail contact forces of an instrumented heavy haul railway vehicle using machine learning}{Mechanical Systems and Signal Processing}
\resumePub{2021}{The effect of railway wheel wear on reprofiling and service life}{Wear}

\vspace{-5pt}
{\scshape Conference Presentations}
\vspace{3pt}

\resumeConf{2026}{International Conference on Railway Technology (Railways 2026)}{Aug. 23--26 (accepted)}{Budapest, Hungary}
\resumeConf{2025}{International Heavy Haul Association (IHHA) Conference}{Nov. 17--21}{Colorado Springs, USA}
\resumeConf{2023}{International Heavy Haul Association (IHHA) Conference}{Aug. 27--31}{Rio de Janeiro, Brazil}
%\resumeConf{2023}{Railway Engineering Symposium}{May 19--20}{Campinas, Brazil}

\vspace{-5pt}


\end{document}